\section{Hyperparameter selection}

Now that the models have been defined in the previous section, an Hyperparameter (HP) selection phase is needed.
The HP choice is essential to obtain the best performance of a neural network, and it can significantly influence the model's ability to generalize to unseen data. 

Only 3 Hyperparameters [with related possible values] will be considered in the models:
\begin{enumerate}
    \item Learning rate: the step size at each iteration while moving toward a minimum of the loss function; it controls how much the model weights are updated during training. [1e-4,1e-1]
    \item Batch Size: the number of samples used to compute each update of the model weights; it influences the stability and speed of the training process. \{2,4,8,16,32\}
    \item Dropout percentage: the fraction of input units to drop during training to prevent overfitting; it helps improve the model's ability to generalize. [0,1]
\end{enumerate}

Given the use of state-of-the-art architectures, other HPs such as activation functions, layers number, and layer types have not been changed.

\subsection{Search Algorithms}
The search for optimal HP values can be done in two ways: manually setting them based on prior experience or 
intuition, or using automated search algorithms.

Although manual testing remains one of the most widely used methods, especially in the student environment, it
requires a deep knowledge of the used algorithms and their HPs meaining. In addition, it has been found that 
a manual procedure is ineffective for several problems due to different factors, such as the high dimensionality 
of the hyperparameter space, the presence of complex interactions between hyperparameters, and the computational 
cost of exhaustively evaluating all possible combinations \cite{YANG2020295}.

Formally, a HPs optimization problem aims to find the HPs configuration $\bm{x}^* \in X$ such that:
$$ \bm{x}^* = \text{arg}\underset{\bm{x}\in X}{\min}f(\bm{x}) $$
where $X$ is the search space containing all the possible hyperparameter configurations, and $f$ is the
objective function (e.g., validation loss or error) to be minimized. A search algorithm is an heuristic algorithm
whose goal is to achieve the optimal, or near-optimal, HPs configuration within a given budget, which can be of
time or computational resources.

\subsubsection{Grid search}
Grid search is the most straightforward hyperparameter tuning algorithm. It discretizes the domain of each hyperparameter and evaluates all possible combinations using cross-validation. The combination that yields the best average performance in cross-validation is selected as the optimal set of hyperparameters.